% Unnumbered chapter (no large chapter number in the heading), in TOC. Do not
% \refstepcounter{chapter} so the following chapter keeps the next number (no
% gap in the sequence). book's \chapter starts with \cleardoublepage/\clearpage;
% relax both so the heading can follow the previous text on the same page.
\begingroup
\let\cleardoublepage\relax
\let\clearpage\relax
\vspace{25pt}
\chapter*{Bicing as a case study}
\phantomsection
\label{ch:use_case_bicing}
\addcontentsline{toc}{chapter}{\protect\numberline{}Bicing as a case study}
\endgroup
\vspace{-45pt} % remove the extra space added by the chapter title

% \chapter* does not advance the chapter counter, so
% floats would otherwise inherit the previous chapter's number. Use Part IV prefix.
\begingroup
\setcounter{figure}{0}
\setcounter{table}{0}
\renewcommand{\thefigure}{IV.\arabic{figure}}
\renewcommand{\thetable}{IV.\arabic{table}}
% Separate IV-* entries from the preceding chapter's figures in the LOF/LOT: this block
% still shares the same internal chapter counter as Chapter~\ref{ch:station_planning_in_barcelona}.
\addtocontents{lof}{\protect\vspace{10pt}}
\addtocontents{lot}{\protect\vspace{10pt}}

Part~\ref{pt:predictive_operations} develops predictive methods for day-to-day \ac{bss} operations, focusing on mobility and battery dynamics as well as component failures. Before presenting those contributions, this section introduces \textit{Bicing} as the empirical setting in which they are developed and tested, outlining how the system has evolved since its launch and how it operates today.

Barcelona’s public docked \ac{bss}, known as \textit{Bicing}, was launched in March 2007 to support the city’s sustainable mobility goals and encourage shifts away from private motorized transport~\cite{Winslow2019}. Implemented with an initial fleet of around 1,500 bicycles and 94 stations, and operated by \textit{Clear Channel}, the system quickly became one of the first large-scale \acp{bss} in Europe. Its rapid success was evident from the fact that, by the end of its first year, \textit{Bicing} had already reached 80,000 subscribers, which was more than six times the number originally anticipated~\cite{Winkels2007}. Moreover, the system’s influence extended beyond its own operations and contributed to a broader cultural shift toward cycling in the city, as reflected in an 82\% increase in bicycle trips in 2007 compared with the previous year~\cite{Winslow2019}. Following its successful debut, \textit{Bicing} expanded progressively over the subsequent years. The system grew outward from the initial downtown stations to cover most flat neighborhoods, reaching 420 docking stations and 6,000 bicycles by 2014~\cite{Winslow2019}. 

A major turning point came in late 2018, when the City Council awarded a new ten-year operating contract to \textit{Pedalem Barcelona}—a joint venture formed by the municipal services company \textit{B:SM (Barcelona de Serveis Municipals)}, \textit{Ferrovial-Cespa}, and \textit{PBSC Urban Solutions}~\cite{ElPais2017}. Under this new contract, the system underwent a comprehensive overhaul in 2019. Key changes included an updated tariff structure, upgraded docking stations, and the introduction of a feature enabling users to reserve bicycles a few minutes in advance. The fleet was also renewed and expanded to around 6,700 bicycles, with the introduction of \acp{e-b} that were capable of pedal assistance up to 25~km/h, making up about 15\% of the bikes initially. Accordingly, stations were modified so that both \acp{e-b} and \acp{m-b} could be accommodated. Moreover, in the following years, the system continued to expand and electrify its fleet. By December 2022, \acp{e-b} made up 47\% of a nearly 7,000-bike fleet. By 2025, the total fleet rose to 8,000 bicycles, with 62\% being \acp{e-b}. The network also increased in coverage, now comprising 557 stations (Figure~\ref{fig:bicing_stations_map}).

\begin{figure}[!hpt]
    \centering
    \includegraphics[width=0.8\textwidth]{0_plots/bicing/bicing_stations_map.jpg}
    \caption[Map of \textit{Bicing} stations]{
        \textbf{\textit{Bicing} stations in 2025.} The black lines delineate the districts of Barcelona.
    }
    \label{fig:bicing_stations_map}
\end{figure}


In its current form, \textit{Bicing} operates 24 hours a day, seven days a week, as a classic docked \ac{bss} (Figure~\ref{fig:bicing_bike_and_station}). Each station consists of a row of docking points where bicycles are locked and can be picked up or returned. Individual stations accommodate between 12 and 54 docking points, and high-demand locations may include multiple stations in close proximity. Since 2019, all docking stations also serve as charging hubs for \acp{e-b}, enabling batteries to recharge automatically when docked.

\begin{figure}[!hpt]
    \centering
    \begin{minipage}[b]{0.48\textwidth}
        \centering
        \includegraphics[width=\textwidth]{0_plots/bicing/bike.jpg}
    \end{minipage}
    \hfill
    \begin{minipage}[b]{0.48\textwidth}
        \centering
        \includegraphics[width=\textwidth]{0_plots/bicing/station.jpg}
    \end{minipage}
    \caption[\textit{Bicing} stations and bikes]{
        \textbf{\textit{Bicing} bikes and stations.} Images obtained from the official \textit{Bicing} website~\cite{BicingOfficial} (left) and the journal \textit{El Jardí}~\cite{ElJardi2024} (right).
    }
    \label{fig:bicing_bike_and_station}
\end{figure}

To use these bicycles, users must be registered members, and casual rentals for tourists are not available by design~\cite{Winslow2019}. This policy is intended to protect local bike-rental businesses and ensure that the system primarily serves residents rather than visitors. Upon subscribing, users receive a contactless card and access credentials for the city’s \textit{SMOU} mobility app, which allows them to unlock and lock bicycles, locate nearby stations, check real-time availability, and reserve a bike up to five minutes in advance. 

Membership is annual, with fees ranging from 35 to 50\texteuro{} depending on the chosen plan. The 50\texteuro{} option offers a flat-rate membership that includes unlimited 30-minute trips on \acp{m-b}, while the 35\texteuro{} option is designed for occasional riders through a pay-per-use model. Both plans apply incremental time-based fees for longer rentals to incentivize short trips. In addition, \acp{e-b} carry a small per-ride cost even within the first 30 minutes.

% Extra vertical space in LOF/LOT before the next numbered chapter's floats: custom
% \thefigure/\thetable do not advance \thechapter, so the usual chapter-change gap is missing.
\addtocontents{lof}{\protect\vspace{10pt}}
\addtocontents{lot}{\protect\vspace{10pt}}

\endgroup








